\subsection{公理}

我们已经看到，特定符号决定了理论 \(\mathfrak{G}\) 中的项和关系。为了构造 \(\mathfrak{G}\) ，我们按以下步骤进行：

\begin{enumerate}
\def\labelenumi{(\arabic{enumi})}
\item
  首先我们在 \(\mathfrak{G}\) 中写下一定数量的关系；这些关系称为 \(\mathfrak{G}\) 的显式公理。出现在显式公理中的字母称为 \(\mathfrak{G}\) 的常量.
\item
  我们规定一条或多条规则\footnote{为简洁起见，这些规则使用§1第1号中提到的符号（尤其是粗斜体字母）来表达；但在规则的表述中完全避免使用这些符号也是容易做到的（见§3第1号，第28页脚注(*)）。}，称为 \(\mathfrak{C}\) 的方案，它们必须具有以下性质：（a）应用这样的规则 \(\mathcal{R}\) 给出 \(\mathfrak{C}\) 中的一个关系； (b) 如果 \(T\) 是 \(\mathfrak{C}\) 中的一项，如果 \(x\) 是一个字母，且如果 \(R\) 是 \(\mathfrak{C}\) 中的一个关系（由方案 \(\mathcal{R}\) 构造），则关系 \((T|x)R\) 也可以通过应用 \(\mathcal{R}\) 得到。
\end{enumerate}

在我们设想的所有情形中，这些条件的验证总是容易的。

通过应用 \(\mathfrak{G}\) 的某个方案构造的每个关系，称为 \(\mathfrak{G}\) 的隐式公理。

直观上，公理要么表示自明的主张，要么表示人们希望从中推出结论的假设。常量表示已明确定义的对象，显式公理所表达的性质被认为对这些对象成立。另一方面，如果字母 x 不是常量，则它表示一个完全未确定的对象；如果通过公理假定对象 x 具有某种性质，则该公理必然是隐式的，因此该性质对任何对象 T 仍然成立。

